% Written By Enoch Yu
% Downloaded from archive.enochyu.com

\documentclass{article}

\usepackage{amsmath}
\usepackage{amsthm}
\usepackage{amssymb}
\usepackage{gensymb}
\usepackage[margin=1in]{geometry}

\usepackage{tikz}
\usetikzlibrary{calc, angles, quotes}
\usepackage{graphicx}
\usepackage{float}
\usepackage[font=small,labelfont=bf]{caption}

\usepackage{parskip}
\usepackage{fancyhdr}
\pagestyle{fancy}
\fancyhead[R]{Enoch Yu}
\pagenumbering{gobble}

\usepackage{enumitem}
\newtheorem{theorem}{Theorem}[section]
\newtheorem{lemma}[theorem]{Lemma}
\newtheorem{sublemma}{Lemma}[section]
\newtheorem{proposition}{Proposition}
\newtheorem{corollary}{Corollary}[theorem]
\newenvironment{solution}{\begin{trivlist}\item[]
{\bf Solution}}{\qed \end{trivlist}}

\pdfinfo{
  /Title (Problem Set 6)
}

\title{Problem Set 6}
\author{Enoch Yu}
\date{May 2025}

\begin{document}

\section*{Probability and Counting}
\[
\boxed{\text{Probability}=
\frac{ \text{cases \textbf{with} order} }
{ \text{total cases \textbf{with} order} }
=\frac{ \text{cases \textbf{without} order} }
{ \text{total cases \textbf{without} order}} }
\]
\subsection*{Example I} There are three red coins
and two blue coins. Find the probability of picking
two coins that are the same color.

\textbf{Solution I: With Order}

Number of cases for choosing coins with same color:
$3 \cdot 2 + 2 \cdot 1$

Total cases: $5 \cdot 4$

\[
\therefore \frac{3 \cdot 2 + 2 \cdot 1}{5 \cdot 4}
= \frac{8}{20} = \boxed{\frac{2}{5}}
\]

\textbf{Solution II: Without Order}

Number of cases for choosing coins with same color:
${}_3C_2 + {}_2C_2$

Total cases: ${}_5C_2$
\[
\therefore \frac{{}_3C_2+{}_2C_2}{{}_5C_2}
= \frac{4}{10} = \boxed{\frac{2}{5}}
\]

To choose the best method to use to solve a problem
depends on each problem. Counting with order may or
may not be more convenient compared to counting without
order.

\subsection*{Example II} Five coins are tossed in the
air simultaneously. What is the probability that there
are two coins that landed on head?

\textbf{Solution I: With Order}

Number of cases of two coins landing on the head:
$\frac{5!}{2!3!}$ (Arrangement of two heads and
three tails)

Total cases: $2^5$
\[
\therefore \frac{ \frac{5!}{2!3!} }{2^5}
= \frac{10}{32} = \boxed{\frac{5}{16}}
\]

\textbf{Solution II: Without Order}

Number of cases of two coins landing on the head:
${}_5C_2$

Total cases:
${}_5C_0 + {}_5C_1 + {}_5C_2 + {}_5C_3 + {}_5C_4 + {}_5C_5$
\[
\therefore \frac{{}_5C_2}
{{}_5C_0 + {}_5C_1 + {}_5C_2 + {}_5C_3 + {}_5C_4 + {}_5C_5}
= \frac{10}{2^5} = \boxed{\frac{5}{16}}
\]

In a more advanced problem, counting with order was more
simple than counting without order.

\newpage
\section*{Problem}
Your wardrobe contains two red socks, two green socks,
two blue socks, and two yello socks. It is currently dark
right now, but you decide to pair up the socks randomly.
What is the probability that none of the pairs are of the
same color?

\textbf{Key Word} Counting Strategy,
Principle of Inclusion and Exclusion

\textbf{Solution I: Without Order}

Total case without order: $\frac{8!}{4!\cdot2^4}$

Principle of Inclusion and Exclusion is applied for
cases without order!!
\[
\frac{8!}{4! \cdot 2^4} - \frac{6!}{3! \cdot 2^3}
\cdot {}_4C_1 + \frac{4!}{2!\cdot2^2} \cdot {}_4C_2
- \frac{2!}{1! \cdot 2^1} \cdot {}_4C_3 +
\frac{0!}{0! \cdot 2^0} \cdot {}_4C_4 = 60
\]
Therefore, $\frac{60}{105} = \boxed{\frac{4}{7}}$
\qed


\textbf{Solution II: With Order}

Let us count by letting that the number of cases does
change when the arrangement of each pair change while
the number remains the same if the arrangement occurred
within a pair.

Total possibility:
${}_8C_2 \cdot {}_6C_2 \cdot {}_4C_2 \cdot {}_2C_2$

For specific case, choose two groups from 4 types of
socks, or ${}_4C_2$. Within each chosen pair, choose
one socks each, or ${}_2C_1 \cdot {}_2C_1$. WLOG, the
next three pairs may be chosen.

Case I: $(1, 1)$

${}_2C_1 \cdot {}_2C_1 = 4$


Case II: $(2,1)$

${}_2C_1 \cdot {}_2C_1 \cdot {}_2C_1
\cdot {}_2C_1 \cdot {}_2C_1 = 32$


Case III: $(2,2)$

${}_2C_1 \cdot {}_2C_1 \cdot {}_4C_2 = 24$


$\therefore \frac{ {}_4C_2 \cdot {}_2C_1 \cdot
{}_2C_1 (4 + 32 + 24) }{ {}_8C_2 \cdot {}_6C_2 
\cdot {}_4C_2 \cdot {}_2C_2 } = \boxed{\frac{4}{7}}$
\qed


If the number of cases increase, PIE, or Solution I,
must be utilized.

\newpage
\section*{Problem}
Given that $468751=5^8+5^7+1$ is a product of two
primes, find both of them.
\begin{solution}
Let $x=5$
\begin{align*}
    x^8 + x^7 + 1
    &= x^8 + x^7 + x^6 - x^6 + 1
    = x^6 (x^2 + x + 1) - x^6 + 1 \\
    &= x^6 (x^2 + x + 1) - x^6 - x^5 - x^4 + x^5 + x^4 + 1
        = x^6 (x^2 + x + 1) - x^4 (x^2 + x + 1)
        + x^5 + x^4 + 1 \\
    &= x^6 (x^2 + x + 1) - x^4 (x^2 + x + 1)
        + x^5 + x^4 + x^3 - x^3 + 1 = x^6 (x^2 + x + 1) \\
    &\qquad \qquad \qquad \qquad \qquad \qquad \qquad
     \qquad \qquad \qquad \qquad
     \text{ } - x^4 (x^2 + x + 1)
     + x^3 (x^2 + x + 1) - x^3 + 1 \\
    &= x^6 (x^2 + x + 1) - x^4 (x^2 + x + 1)
        + x^3 (x^2 + x + 1) - (x - 1)(x^2 + x + 1) \\
    &= (x^2 + x + 1)(x^6 - x^4 + x^3 - x + 1) \\
    &= 31 \cdot 15121
\end{align*}
Therefore, the two prime factors are
$\boxed{31 \text{ and } 15121}$.
\end{solution}
\end{document}
